\chapter{Comparison and generalized rules}
\label{chap:comparison-and-rules}

This chapter begins the internal-proof layer.  It consumes the explicit literature and computation inputs above to prove comparison, page-extension, Leibniz, Moss-adapter, and Mahowald results inside the project.

\section{Extension-SS consumers}

% SOURCE: AIM-7 concrete tracer bullet; MainPaper/main.tex:430-453.
\begin{proposition}[Concrete eta detection slice]
  \label{prop:classical-eta-extension-iff-detected}
  \uses{def:classical-eta-ess-concrete,prop:eta-ess-regression}
  \lean{KIP126.Classical.ExtensionSS.extension_iff_detected}
  \lean{KIP126.Classical.ExtensionSS.FExtension}
  \lean{KIP126.Classical.ExtensionSS.DetectedBy}
  \lean{KIP126.Classical.ExtensionSS.Essential}
  \lean{KIP126.Classical.ExtensionSS.Crossing}
  \notready
  For the concrete finite eta rows, extension is equivalent to membership in
  the set-valued detection predicate.  Essentiality and crossing are defined
  only on this construction, so detection representatives are not forced to
  be unique.
\end{proposition}

\begin{proof}
  Expand the two predicates as membership in the detected and differential
  sets and rewrite with the concrete input's equality between those sets.
  The eta rows themselves are supplied by the catalogued regression evidence.
\end{proof}

% SOURCE: AIM-7 concrete tracer bullet; MainPaper/main.tex:430-453.
\begin{definition}[Concrete classical $\eta$-ESS tracer bullet]
  \label{def:classical-eta-ess-concrete}
  \uses{def:classical-adams-ss,def:ess-complex-shape,thm:external-adams-h4-d2-slice,prop:eta-ess-regression,def:catalogued-external-evidence}
  \lean{KIP126.Classical.ExtensionSS.etaESS}
  \lean{KIP126.Classical.ExtensionSS.E₀}
  \lean{KIP126.Classical.ExtensionSS.abutment}
  \lean{KIP126.Classical.ExtensionSS.ClassicalEtaESSAdapter}
  \notready
  On the AIM-5 concrete classical Adams data, the $\eta$ tracer bullet uses
  Mathlib's `CategoryTheory.SpectralSequence` with page-zero object
  $E_\infty(X)\oplus E_\infty(Y)$, differential degree $(n,n)$, and
  componentwise abutment $\ker(\pi_*\eta)\oplus\operatorname{coker}(\pi_*\eta)$.
  The finite differential rows carry the AIM paper locator and are injected
  through the `etaEssRegression` claim ledger; this does not add a global
  differential axiom.
\end{definition}

% SOURCE: MainPaper/main.tex:327-342, unlabelled Corollary 2.6.
\begin{corollary}[Map filtration forces initial ESS differentials to vanish]
  \label{cor:fess-below-map-filtration-zero}
  \uses{thm:external-map-filtration-factorization,prop:fextension-iff-detected}
  \notready
  If $\operatorname{AF}(f)=k$, then $d_i^f=0$ for every $i<k$.
\end{corollary}

\begin{proof}
  Factor $f$ into $k$ homology-trivial maps.  Each factor raises Adams
  filtration by at least one, so the composite raises it by at least $k$.
  Proposition~\ref{prop:fextension-iff-detected} then excludes a nonzero
  target in any smaller filtration shift.
\end{proof}

\section{Moss convergence adapter}

\begin{proposition}[Mapping-stage composition at a left successor]
  \label{prop:moss-stage-composition-successor}
  \uses{def:hf2-adams-tower,prop:h6-fiber-smash-step,prop:h6-tower-smash-comparison,prop:h6-smash-pairing-left-transition}
  \lean{KIP126.Classical.Adams.Moss.stageComposition_succ_comparison}
  \leanok
  Fix an arbitrary unit $u:\mathbb 1\to H$ in a braided monoidal closed
  stable category.  Assume every right tensor functor is triangulated with
  its chosen suspension comparison.  Write $T_s[X,Y]$ for the actual
  mapping-object Adams stage, $\overline H=\operatorname{fib}(u)$, and
  \[
    \chi_s:T_{s+1}[X,Y]\xrightarrow{\ \cong\ }
      \overline H\otimes T_s[X,Y]
  \]
  for the existing fiber comparison.  If $c_{s,t}$ denotes the fixed
  stage composition with output $T_{t+s}[X,Z]$, then
  \[
    \chi_{t+s}\circ c_{s+1,t}
      =(1_{\overline H}\otimes c_{s,t})\circ\alpha
        \circ(\chi_s\otimes1_{T_t[Y,Z]}).
  \]
  This identifies the actual successor map before any fiber inclusion is
  applied.  It requires no coefficient multiplication or boundary input.
  It does not prove that an ordered right-boundary defect vanishes, nor
  provide higher-page Leibniz or boundary compatibility.
\end{proposition}
\begin{proof}
  \leanok
  In the iterated-smash coordinates, separate the additional left factor
  $\overline H$.  Associator naturality and the braided interchange
  identity express successor composition as
  $\overline H\otimes c_{s,t}$ after reassociation.  Transport this
  equality through the actual tower-smash comparisons and cancel the
  output comparison isomorphism.  The successor definition of the fiber
  comparison gives exactly the displayed equality.
\end{proof}

\begin{proposition}[Actual long-layer composition obstruction]
  \label{prop:moss-long-layer-composition-obstruction}
  \uses{def:hf2-adams-tower,def:h6-tower-long-layer,thm:stable-homotopy-long-exact}
  \lean{KIP126.Classical.Adams.Moss.longLayerProjectedComposition_ι,KIP126.Classical.Adams.Moss.longLayerCompositionBoundary_ι,KIP126.Classical.Adams.Moss.longLayerCompositionBoundary_rightTower_factors,KIP126.Classical.Adams.Moss.longLayerCompositionBoundary_leftTower_factors,KIP126.Classical.Adams.Moss.longLayerTwoCompositionObstruction,KIP126.Classical.Adams.Moss.longLayerTwoComposition_exists_iff_obstruction_eq_zero,KIP126.Classical.Adams.Moss.longLayerTwoCompositionObstruction_ι}
  \leanok
  Use the actual Adams towers of the three mapping objects, the specified
  coefficient multiplication, and the selected triangulated left and right
  tensor functors with their suspension comparisons.  For any length
  $r\ge1$ and nonnegative tower levels $s,t$,
  the fixed projected long-layer composition agrees on two tower images
  with stage composition followed by the output layer inclusion.  Its
  connecting map therefore vanishes on those two tower images.
  With either input restricted to its tower, exactness factors this boundary
  through the other input's actual long-cofiber connecting map, retaining
  the tensor suspension comparison.

  At length two, write $B$ for the fixed composition boundary and put
  $O=(\Sigma\iota_{s+t+1})\circ B$, where $\iota_{s+t+1}$ is the actual
  next tower-to-layer map.  A projection-compatible length-two composition
  exists if and only if $O=0$.  The obstruction vanishes on two tower
  images.  This restricted vanishing alone does not establish $O=0$ on
  the entire tensor of long layers.  Global length-two vanishing under
  the additional coefficient-unit compatibility hypotheses is proved below.
  The two restricted factorizations do not supply a full two-term boundary
  formula or the pagewise Leibniz and convergence compatibility needed by Moss.
\end{proposition}
\begin{proof}
  \leanok
  The coefficient-unit identity and the actual layer projection identify
  the composite on the two tower images with stage composition followed by
  the output inclusion.  Consecutive maps in its cofiber triangle compose
  to zero, giving the restricted boundary vanishing.  Tensor the left
  long-cofiber triangle with the right tower and apply Yoneda exactness to
  factor the boundary restricted on the right.  Apply the same argument
  to the right long-cofiber triangle tensored with the left tower for the
  other factorization; keep both selected suspension comparisons.

  Exactness for the output long-cofiber triangle first expresses a
  projection-compatible composition as a lift of its fixed boundary.
  At length two, apply exactness to the shifted next-layer triangle:
  such a lift exists precisely when composition with the shifted inclusion
  vanishes.  This composite is $O$; the shift signs are absorbed by negating
  the chosen lift.  Its vanishing on two tower images follows by composing
  the already proved restricted boundary equality with that same inclusion.
\end{proof}

\begin{proposition}[First-boundary cycles as a unit equalizer]
  \label{prop:moss-first-boundary-unit-equalizer}
  \uses{def:hf2-adams-tower,def:h6-ring-cooperations,def:h6-coefficient-layer-product,thm:stable-homotopy-long-exact}
  \lean{KIP126.Classical.Adams.Moss.coefficientPairing_precomp,KIP126.Classical.Adams.Moss.coefficientPairing_postcomp,KIP126.Classical.Adams.Moss.coefficientPairing_innerUnit,KIP126.Classical.Adams.Moss.coefficientPairing_preserves_unit_equalizer,KIP126.Classical.Adams.Moss.firstBoundary_eq_zero_iff_unit_equalizer,KIP126.Classical.Adams.Moss.coefficientPairing_firstBoundary_eq_zero,KIP126.Classical.Adams.Moss.layerFirstBoundary_comparison,KIP126.Classical.Adams.Moss.layerFirstBoundary_eq_zero_iff_unit_equalizer,KIP126.Classical.Adams.Moss.layerComposition_firstBoundary_eq_zero}
  \leanok
  Fix the actual coefficient object $H$, its unit $\eta$, and a specified
  coefficient ring structure.  Assume $H\otimes-$ is triangulated with
  its chosen suspension comparison, and the natural transformation
  $\eta:\operatorname{Id}\to H\otimes-$ commutes with suspension.
  Write $I A=\operatorname{fib}(\eta_A)$ and let
  \[
    \partial_A:H\otimes A\xrightarrow{\delta_A}\Sigma I A
      \xrightarrow{\Sigma\eta_{IA}}\Sigma(H\otimes I A)
  \]
  be the actual first resolution boundary.  For every test object $U$ and
  every morphism $a:U\to H\otimes A$,
  \[
    \partial_A a=0\quad\Longleftrightarrow\quad
    \eta_{H\otimes A}a=(1_H\otimes\eta_A)a.
  \]
  In a braided category, the specified coefficient pairing associated to
  any $f:A\otimes B\to D$ preserves this equalizer, hence sends two
  first-boundary cycles to a first-boundary cycle.

  The actual Adams layer comparisons transport the same criterion to
  each nonnegative tower level.  For mapping objects, additionally assume
  the category is monoidal closed and every right tensor functor is
  triangulated with its chosen suspension comparison.  Then the fixed
  layer composition preserves these kernels on morphisms from arbitrary
  test objects; no assertion that homotopy groups detect maps is required.
\end{proposition}
\begin{proof}
  \leanok
  Unit naturality identifies the first boundary with the connecting map
  in the triangle obtained by tensoring the unit fiber triangle with $H$.
  Exactness factors a zero-boundary morphism through the inner unit
  insertion.  The coefficient action retracts both inner and outer unit
  insertion, forcing the factor to be the original morphism $a$.
  Conversely, the equalizer identity reduces the boundary to a composite
  of consecutive maps in that distinguished triangle.
  Naturality of the actual coefficient pairing and its two unit identities
  preserve the equalizer.  Finally, use the specified layer isomorphisms
  and the comparison of their connecting maps to transport this result
  to the actual layer composition, retaining the shift signs.
\end{proof}

\begin{proposition}[Actual composition on length-two long layers]
  \label{prop:moss-long-layer-two-composition}
  \uses{prop:moss-long-layer-composition-obstruction,prop:moss-first-boundary-unit-equalizer}
  \lean{KIP126.Classical.Adams.Moss.longLayerTwoProjection_firstBoundary_eq_zero,KIP126.Classical.Adams.Moss.longLayerTwoCompositionObstruction_eq_zero,KIP126.Classical.Adams.Moss.longLayerTwoComposition_exists}
  \leanok
  In a braided monoidal closed stable category, fix the actual coefficient
  ring structure, triangulated right tensor functors with suspension
  comparisons, a triangulated $H\otimes-$ with its suspension comparison,
  and suspension compatibility of the coefficient-unit natural
  transformation.  For any mapping objects $[X,Y],[Y,Z],[X,Z]$ and
  nonnegative levels $s,t$, the actual obstruction $O$ from
  Proposition~\ref{prop:moss-long-layer-composition-obstruction} vanishes
  on the entire tensor of length-two long layers.  Consequently, writing
  $L_2^s$ for those long layers, there is a morphism
  \[
    \mu:L_2^s[X,Y]\otimes L_2^t[Y,Z]\longrightarrow L_2^{s+t}[X,Z]
  \]
  whose composite with the output layer projection is exactly the fixed
  projected composition.
  This proves existence at length two under the stated coherence
  hypotheses.  It does not choose compatible compositions for every
  length, prove a pagewise Leibniz rule, or establish Moss convergence.
\end{proposition}
\begin{proof}
  \leanok
  Each length-two projection is an actual first-boundary cycle, by its
  cofiber-factorization diagram and the zero composite of consecutive
  triangle maps.  The unit-equalizer result shows that their layer product
  is also a first-boundary cycle.  Its first boundary, after the actual
  level reindexing, is exactly $O$, so $O=0$ globally.  Apply the proved
  length-two obstruction criterion to obtain the required lift.
\end{proof}

\begin{proposition}[Actual one-sided boundary and second-cycle composition]
  \label{prop:moss-boundary-and-two-cycles}
  \uses{prop:moss-stage-composition-successor,prop:moss-long-layer-two-composition}
  \lean{KIP126.Classical.Adams.Moss.stageComposition_δ_left,KIP126.Classical.Adams.Moss.layerComposition_ι_right_comparison,KIP126.Classical.Adams.Moss.layerComposition_ι_right_δ_comparison,KIP126.Classical.Adams.Moss.layerComposition_ι_right_δ,KIP126.Classical.Adams.Moss.longLayerTwoComposition,KIP126.Classical.Adams.Moss.longLayerTwoPairing,KIP126.Classical.Adams.Moss.longLayerTwoComposition_projection,KIP126.Classical.Adams.Moss.longLayerTwoPairing_projection,KIP126.Classical.Adams.Moss.firstComposition_mem_cycles_two,KIP126.Classical.Adams.Moss.twoCycleComposition}
  \leanok
  In the braided monoidal closed stable category, assume right tensor
  functors are triangulated with suspension comparisons, and supply their
  explicit \texttt{RightTensorSuspensionCompatibility}.  The actual
  successor composition intertwines the unit-fiber connecting maps.
  For the specified coefficient ring structure and monoidal preadditivity,
  this gives the actual layer identity
  \[
    \delta_{t+s}\circ P_{s,t}\circ(1_{L_s}\otimes q_t)
      =(\Sigma c_{s+1,t})\circ\sigma\circ(\delta_s\otimes1_{T_t}),
  \]
  where $q_t:T_t\to L_t$ is the tower-to-layer map and $\sigma$ is the
  chosen right-tensor suspension comparison.  The right input is restricted
  to its tower image; this is one boundary identity, not a full Leibniz law.

  Separately, under the hypotheses of
  Proposition~\ref{prop:moss-long-layer-two-composition} and monoidal
  preadditivity, choose a proved length-two composition lift.  Its induced
  homotopy pairing commutes with the actual long-layer projections, so the
  specified first-layer composition restricts to a bilinear map
  \[
    Z_2^{s,u}[X,Y]\times Z_2^{t,v}[Y,Z]
      \longrightarrow Z_2^{s+t,u+v}[X,Z].
  \]
  These are the actual second-page cycle modules, before quotienting by
  boundaries.  No boundary-ideal condition, $E_2$ descent, compatibility of
  the chosen lifts across lengths, or Moss convergence follows here.
\end{proposition}
\begin{proof}
  \leanok
  The successor comparison and unit-fiber suspension compatibility give
  the first connecting-map identity.  Transport through the actual layer
  isomorphisms; their two boundary signs cancel in the displayed formula.
  For the second assertion, select the existing length-two lift and use
  naturality of the homotopy tensor pairing to prove projection
  compatibility.  The characterization of cycles as images of long-layer
  projections then supplies cycle preservation and the bilinear restriction.
\end{proof}

% SOURCE: MainPaper/main.tex:2537--2545; Moss Theorem 1.2.
% The crossing and weak-convergence conventions have also been checked against
% Belmont--Kong, arXiv:2112.08689v2, Definitions 2.3--2.4.
\begin{theorem}[Moss convergence adapter]
  \label{thm:moss-convergence-adapter}
  \uses{def:toda-bracket,def:massey-product,thm:toda-product-identities,def:adams-detection,def:external-result,prop:moss-long-layer-composition-obstruction}
  \lean{KIP126.Classical.Adams.Moss.PageMassey.Relation,KIP126.Classical.Adams.Moss.PageMassey.Indeterminacy,KIP126.Classical.Adams.Moss.StatementAt,KIP126.Challenge2.MossInterface,KIP126.Challenge2.StandardSphereMossInterface,KIP126.Interface.Solution.sphereMossInterface}
  \notready
  Fix one composition system whose page values are constrained by the actual
  long-layer representatives, with Leibniz, associativity and next-page
  compatibility, and compatible detection in the actual tower-image filtration.
  For $r\ge3$, the Massey relation on $E_r$ uses a defining system on
  $E_{r-1}$: $dy=ab$, $dz=bc$, and the class of
  $yc-(-1)^{|a|}az$. Its indeterminacy is the sum of the two corresponding
  composition images on $E_r$.

  Suppose $a,b,c$ detect specified stable maps, both products vanish on the
  page and as stable composites, and the two product mapping towers satisfy
  weak convergence and the Moss-direction noncrossing conditions. Then some
  Massey class is permanent and detects an element of the Toda bracket.
  Neither zero indeterminacy nor survival of every Massey class is assumed.
  The chosen family controls the convergence range; no convergence of every
  ambient object is asserted.

  The precise interface and its existence obligation for the fixed sphere
  mapping model are now written. The mapping-sphere comparison with the
  previously fixed sphere Adams object, the model data, and the convergence
  proof remain unfinished. This preceding-page definition does not claim to
  define the separate cobar Massey product on $E_2$.
\end{theorem}

\section{Synthetic comparison and normalized lifts}

% SOURCE: MainPaper/main.tex:685-691; AIM-6 concrete synthetic comparison slice.
\begin{proposition}[The $h_4$ classical--synthetic $d_2$ comparison]
  \label{prop:aim6-h4-comparison}
  \uses{def:classical-adams-ss,def:synthetic-adams-ss,def:synthetic-tridegree,def:reindexed-ss-morphism,def:external-result}
  \lean{KIP126.Synthetic.SpectralSequence.SyntheticLambdaAction}
  \lean{KIP126.Synthetic.SpectralSequence.lambdaMapFromAction}
  \lean{KIP126.Synthetic.SpectralSequence.SyntheticLambdaAction.lambdaMap_comm}
  \lean{KIP126.Synthetic.SpectralSequence.weightPreserving_differential}
  \lean{KIP126.Comparison.ClassicalSynthetic.sourcePageMap}
  \lean{KIP126.Comparison.ClassicalSynthetic.sourcePageMap_differential_naturality}
  \notready
  The internal \texttt{SyntheticLambdaAction} preserves the actual cycle
  and boundary towers.  Its page maps are induced quotient maps of the same
  ambient map; differential commutation is an explicit compatibility field.
  Fixed-weight classical--synthetic comparison uses the same internal model
  and the same requirement.  Weight preservation follows from the fixed
  differential degree, without an independently supplied field.

  The $h_4$ correspondence remains \texttt{notready}: the named standard
  classes must be identified on the selected objects, the action must be
  identified with the regraded image of the actual synthetic deformation map,
  and the representative comparison across weights must respect that action.
  Generic synthetic sequences no longer carry arbitrary independently chosen
  $h_4$ and $h_0h_3^2$ elements.
\end{proposition}

\begin{proof}
  The AIM-5 catalogue wrapper supplies the classical statement and degree
  equalities.  The remaining proof obligation is the representative-induced
  comparison across the two fixed weights and its compatibility with the
  typed action; no proof of the differential correspondence is claimed here.
\end{proof}

% SOURCE: MainPaper/main.tex, Theorem thm:17e90ac0(1);
% Source/BHS/source/SynRevBigraded.tex:59-63, lemm:ctauE2;
% Source/BurklundXu/paper.txt:2996-3001, Lemma 7.5.
% The shifted form additionally requires the actual shift/quotient comparison.
\begin{definition}[First-quotient homotopy comparison]
  \label{def:synthetic-first-quotient-homotopy-comparison}
  \uses{def:synthetic-sphere,def:synthetic-lambda-quotient,def:classical-adams-ss}
  \lean{KIP126.Challenge2.FirstQuotientHomotopyComparison}
  \leanok
  For a fixed classical object $X$, coefficient object $H$, and synthetic
  analogue functor $\nu$, a first-quotient homotopy comparison consists of
  additive isomorphisms, for all integers $a,s,t$,
  \[
    \pi_{t-s,t+a}\bigl((\Sigma^{0,a}\nu X)/\lambda\bigr)
       \cong E_2^{s,t}(X).
  \]
  The source is the actual quotient of the shifted object, and the target
  is the actual second page of the same classical Adams tower.
  The unshifted case is the special-fiber homotopy comparison in the
  cited sources; extending it to arbitrary $a$ also requires compatibility
  of the actual shifts, $\lambda$, and cofiber quotients.

  Both sides are already defined, so this is a precise parameterized
  delivery type.  Its model witness has not been constructed.
  The existing special-fiber $E_\infty$ presentation identifies associated
  graded groups and does not by itself provide these homotopy isomorphisms.
  This comparison type does not include naturality or Bockstein-differential
  compatibility conditions.
\end{definition}

% PROOF-SOURCE: BHS Theorem A.1 and its special-fiber edge isomorphism;
% MainPaper/main.tex:813-823.  This is the exponent-one boundary case that is
% not covered by MainPaper/main.tex:849-855.
\begin{proposition}[$E_\infty$ at the first $\lambda$-quotient]
  \label{prop:synthetic-einfty-first-quotient}
  \uses{thm:external-lambda-bockstein,thm:external-synthetic-rigidity,def:synthetic-lambda-quotient,def:synthetic-adams-ss,def:ss-convergence-detection-witness,def:bhs-aim-regrading}
  \notready
  For every admissible $X$, the special-fiber edge map identifies the
  synthetic Adams abutment of $\nu X/\lambda$ with the classical initial page:
  \[
    E_\infty^{s,t,w}(\nu X/\lambda)\cong
    \begin{cases}
      E_2^{s,t}(X),&w=t,\\
      0,&w\ne t.
    \end{cases}
  \]
  Under this isomorphism the edge class has a unique detected homotopy class
  in its tridegree.  This proposition is the explicit exponent-one interface;
  it does not extend the finite $E_\infty$ formula $Z_i/B_j$ of
  Theorem~\ref{thm:external-synthetic-einfty-quotient} outside its stated
  range $r\ge2$.
\end{proposition}

\begin{proof}
  Apply the synthetic Adams $E_2$ functor to the cofiber of $\lambda$.
  Rigidity identifies its page with
  $E_2^{*,*}(X)\otimes\mathbb F_2[\lambda]/(\lambda)$, concentrated in weight
  $w=t$.  The $\lambda$-Bockstein comparison gives the edge isomorphism
  $E_2^{s,t}(X)\cong\pi_{t-s,t}(\nu X/\lambda)$.  Strong convergence and the
  fact that this edge map already exhausts the abutment force every
  differential and hidden additive extension at this quotient to vanish.
  Regrade to AIM's three coordinates to obtain the displayed presentation and
  uniqueness statement.
\end{proof}

% SOURCE: MainPaper/main.tex:865-875, load-bearing compatibility remark after Propositions 3.11-3.12.
\begin{proposition}[Compatibility of $E_\infty$ formulas with $\lambda$ and $\rho$]
  \label{prop:synthetic-einfty-map-compatibility}
  \uses{thm:external-synthetic-einfty-nu,thm:external-synthetic-einfty-quotient,prop:synthetic-einfty-first-quotient,def:synthetic-lambda-quotient,def:bhs-aim-regrading}
  \notready
  For $r'\ge r\ge2$, multiplication is the degree-zero synthetic map
  \[
    \lambda^{r'-r}:
    \Sigma^{0,r-r'}(\nu X/\lambda^r)\longrightarrow\nu X/\lambda^{r'}.
  \]
  Under the preceding isomorphisms, its map on $E_\infty$ is
  \[
    E_\infty^{s,t,w}(\nu X/\lambda^r)
      \longrightarrow
    E_\infty^{s,t,w-r'+r}(\nu X/\lambda^{r'}),
  \]
  and, when the displayed cycle and boundary indices are positive, it is the
  quotient map
  \[
    \frac{Z_{r-t+w}^{s,t}(X)}{B_{1+t-w}^{s,t}(X)}
      \longrightarrow
    \frac{Z_{r-t+w}^{s,t}(X)}{B_{1+t-w+r'-r}^{s,t}(X)}.
  \]
  For $r\ge r''\ge2$, reduction
  $\rho:\nu X/\lambda^r\to\nu X/\lambda^{r''}$
  is, in the same range, the inclusion
  \[
    \frac{Z_{r-t+w}^{s,t}(X)}{B_{1+t-w}^{s,t}(X)}
      \hookrightarrow
    \frac{Z_{r''-t+w}^{s,t}(X)}{B_{1+t-w}^{s,t}(X)}.
  \]
  The project's zero-index convention supplies the remaining cycle- and
  boundary-subscript cases within these exponent ranges.

  For the boundary case $r=1<r'$, the same suspended multiplication map sends
  the special-fiber edge group to the indicated weight of the $r'$ quotient;
  it is the canonical quotient and hence is surjective or zero in each
  tridegree.  For $r''=1<r$, reduction identifies the surviving-cycle subgroup
  with a subgroup of the special-fiber edge group and hence is injective or
  zero.  When both exponents are one the map is the identity.  These boundary
  statements use Proposition~\ref{prop:synthetic-einfty-first-quotient}, not
  the finite-quotient formula outside its range.

  There are also two explicit infinite-endpoint clauses.  For every $n>0$,
  multiplication
  \[
    \lambda^n:\Sigma^{0,-n}\nu X\longrightarrow\nu X
  \]
  is described as follows.  On a tridegree $(s,t,w)$ of the suspended domain,
  put $j=1+t-(w+n)$.  When the source is nonzero, suspension identifies it
  with $Z_\infty^{s,t}(X)/B_j^{s,t}(X)$ and the map is the
  boundary-cutoff quotient
  \[
    Z_\infty^{s,t}(X)/B_j^{s,t}(X)
      \longrightarrow
    Z_\infty^{s,t}(X)/B_{j+n}^{s,t}(X),
  \]
  since the target cutoff is $1+t-w=j+n$.  It is therefore surjective or zero
  in each tridegree.  For reduction
  \[
    \rho:\nu X\longrightarrow\nu X/\lambda^n
  \]
  fix a tridegree $(s,t,w)$ and put $j=1+t-w$.  It is the cycle-cutoff inclusion
  \[
    Z_\infty^{s,t}(X)/B_j^{s,t}(X)
      \hookrightarrow
    Z_{n-t+w}^{s,t}(X)/B_j^{s,t}(X)
  \]
  whenever $n\ge2$ and $0\le t-w<n$, and is injective or zero otherwise.  At
  $n=1$ it is the inclusion of permanent representatives
  into the special-fiber edge group (or zero off its unique weight).  The
  infinite source uses Theorem~\ref{thm:external-synthetic-einfty-nu}; it is
  not obtained by substituting $r=\infty$ into the finite formula.
\end{proposition}

\begin{proof}
  Express both synthetic maps on the $\lambda$-Bockstein exact couple.  The
  first changes the boundary cutoff by $r'-r$ and shifts the weight by
  $r-r'$, while the second changes the cycle cutoff without shifting weight.
  For exponents at least two, transport these maps through the two external
  $E_\infty$ isomorphisms and check the zero-index edge cases separately.  If
  an endpoint exponent is one, use the special-fiber edge presentation: the
  multiplication map quotients the initial-page group by the later boundary
  cutoff, while reduction includes the surviving-cycle subgroup.  This gives
  the asserted surjective/zero and injective/zero alternatives without
  applying the finite-quotient theorem at exponent one.

  For an infinite endpoint, compare the untruncated formula
  $Z_\infty/B_j$ directly with its weight-shifted copy and with the finite
  quotient formula.  Multiplication by $\lambda^n$ raises the boundary cutoff
  from $j$ to $j+n$, hence is the displayed quotient.  Reduction retains the
  boundary cutoff and includes $Z_\infty$ into the finite cycle subgroup; for
  exponent one use the special-fiber edge group.  These calculations prove
  the two infinite clauses independently of any limiting shorthand.
\end{proof}

% SOURCE: MainPaper/main.tex:878-910, Example exam:synEinfty.
\begin{proposition}[Synthetic $14$-stem regression]
  \label{prop:synthetic-14-stem-regression}
  \uses{thm:external-synthetic-rigidity,thm:external-synthetic-einfty-nu,thm:external-synthetic-einfty-quotient,def:external-evidence}
  \notready
  Given the explicit classical evidence
  $d_2(h_4)=h_0h_3^2$,
  $d_3(h_0h_4)=h_0d_0$, and
  $d_3(h_0^2h_4)=h_0^2d_0$, the synthetic spectral sequences of
  $S^{0,0}$, $S^{0,0}/\lambda^2$, and $S^{0,0}/\lambda^3$ have exactly the
  differentials, permanent multiples, and $\lambda$-torsion classes listed in
  Example 3.15.
\end{proposition}

\begin{proof}
  Apply rigidity to insert the factors $\lambda$ and $\lambda^2$, truncate at
  the quotient exponent, and evaluate the two $E_\infty$ formulas.  A finite
  degree check verifies the list and serves as a regression test for grading
  and zero-index conventions.
\end{proof}

% SOURCE: MainPaper/main.tex:946-961, cofiber assertion in Notation not:fhat.
\begin{proposition}[Cofiber of a normalized synthetic map]
  \label{prop:synthetic-hat-cofiber}
  \uses{def:synthetic-hat-map,def:normalized-synthetic-exactness-case,thm:external-nu-cofiber-criterion,thm:external-synthetic-triangle-lift}
  \notready
  Let $X\xrightarrow fY\xrightarrow gC(f)$ be a cofiber sequence equipped with
  a normalized cofiber exactness case.  Let $\widehat f$ and $\widehat g$ be
  the triangle-compatible components obtained simultaneously from the same
  lifted rotation carried by that case.  Then
  \[
    C(\widehat f)\simeq\Sigma^{0,-e(g)}\nu C(f).
  \]
  Thus the tagged zero and monomorphism cases, both of which explicitly have
  $e(g)=0$, give no shift.  The tagged epimorphism case explicitly has
  $e(g)=1$ and gives a $(0,-1)$ shift.  The homology-map adjective alone does
  not select a shift in a degenerate case.  No such equivalence is asserted for
  $\lambda^{k-1}\widetilde f$ formed from an arbitrary full lift.
\end{proposition}

\begin{proof}
  Eliminate the tagged exactness datum.  In the zero constructor, rotate so
  that $f$ is the positive-filtration connecting map; in the monomorphism
  constructor, use the original short exact homology sequence; in the
  epimorphism constructor, rotate so that $g$ is the positive-filtration
  connecting map.
  Apply the cofiber criterion to the other two components in that chosen
  rotation.  Because $\widehat f$ and $\widehat g$ are components of this one
  lifted triangle, its cofiber identifies with the displayed object.  After
  rotating back, use the constructor's explicit $e(g)$ field to obtain exactly
  the asserted shift.  The argument never
  replaces $\widehat f$ by a torsion-equivalent arbitrary full lift.
\end{proof}

% SOURCE: MainPaper/main.tex:965-971, unlabelled Proposition 3.20.
\begin{proposition}[Normalized synthetic triangle]
  \label{prop:normalized-synthetic-triangle}
  \uses{def:synthetic-hat-map,def:normalized-synthetic-exactness-case,prop:positive-adams-filtration-homology-zero,prop:synthetic-hat-cofiber,thm:external-synthetic-triangle-lift,def:cofiber-triangle,def:hf2-homology}
  \notready
  If $X\xrightarrow fY\xrightarrow gZ\xrightarrow h\Sigma X$ is a
  distinguished triangle and $e(f)+e(g)+e(h)=1$, there is a simultaneous
  choice of triangle-compatible normalized lifts that forms a distinguished
  synthetic triangle
  \[
    \nu X\xrightarrow{\widehat f}\Sigma^{0,-e(f)}\nu Y
    \xrightarrow{\widehat g}\Sigma^{0,-e(f)-e(g)}\nu Z
    \xrightarrow{\widehat h}\Sigma^{0,-1}\nu\Sigma X.
  \]
\end{proposition}

\begin{proof}
  The sum-one hypothesis selects exactly one of the three explicit
  $e$-patterns.  Proposition~\ref{prop:positive-adams-filtration-homology-zero}
  makes the unique positive-filtration map zero on homology.  Exactness of the
  cofiber long exact sequence then says respectively that $H_*f$ is zero,
  injective, or surjective, and supplies the corresponding rotated short exact
  sequence.  Transport along the chosen equivalence with the selected cofiber
  triangle, rotate so that this positive-filtration map is the connecting map,
  and apply the triangle-lift theorem once.  These fields form the appropriate
  tagged exactness datum.  Rotate that same synthetic triangle back.  Its three components define the normalized lifts
  simultaneously; Proposition~\ref{prop:synthetic-hat-cofiber} identifies the
  required cofibers.  The equality of the three $e$-values makes the total
  weight shift exactly $-1$.
\end{proof}

\section{Classical--synthetic extension comparison}

% PROOF-SOURCE: the standard obstruction calculation for a countable inverse
% system of abelian torsors, combined with BHS Proposition A.13.
\begin{lemma}[Mittag--Leffler criterion for an inverse-limit ESS extension]
  \label{lem:synthetic-ess-coherent-tower-limit}
  \uses{def:synthetic-ess-solution-tower,prop:synthetic-ess-coherent-tower-of-surjective,prop:synthetic-ess-coherent-tower-of-finite,prop:lambda-quotient-restriction-coherence,thm:external-synthetic-lambda-complete}
  \notready
  In the setting of
  Definition~\ref{def:synthetic-ess-solution-tower}, suppose every finite
  solution fiber is nonempty.  There is a canonical torsor obstruction
  \[
    \omega(\operatorname{Sol}_\bullet)
      \in \varprojlim^{1}_{m} K_m
  \]
  whose vanishing is equivalent to the existence of a coherent solution
  tower.  Any of the following is sufficient for vanishing:
  \begin{enumerate}
    \item every finite solution lifts, that is, each
      $\operatorname{res}_{m+1,m}$ is surjective;
    \item the difference-group tower $(K_m)$ is Mittag--Leffler: for every
      $m$, the images of $K_l\to K_m$ stabilize as $l$ increases.  Surjective
      transition maps are a sufficient special case.  Then
      $\varprojlim^1_m K_m=0$.
    \item every actual strict solution fiber is finite.  Together with
      nonemptiness, Proposition~\ref{prop:synthetic-ess-coherent-tower-of-finite}
      constructs some coherent tower, without fixing its initial value.
  \end{enumerate}
  For an untruncated conclusion, additionally require the following actual
  representative comparison.  Let the maps $F_m$ be reductions of a map
  $F_\infty$ between $\lambda$-complete synthetic spectra, and let
  $R_m,E_m,\Phi_m$ be the representative groups, equation groups, and
  equation maps defining their strict fibers.  Require the canonical
  restriction maps to induce isomorphisms
  \[
    c_R:R_\infty\xrightarrow{\ \cong\ }\varprojlim_m R_m,
    \qquad
    c_E:E_\infty\xrightarrow{\ \cong\ }\varprojlim_m E_m,
    \qquad
    c_E\Phi_\infty=(\varprojlim_m\Phi_m)c_R,
  \]
  carrying the actual untruncated equation-label triple to the compatible
  finite equation-label triples.  These are comparisons of the actual
  filtered source, target, and map, including their fixed cycle labels.
  Under these additional hypotheses, a coherent solution tower determines
  the untruncated relation
  \[
    d_n^{F_\infty}(x_\infty)=y_\infty
  \]
  in the complete target coset, and every untruncated representative
  solution restricts to a coherent finite tower.  Constructing these
  comparisons from the actual $\lambda$-complete synthetic model remains
  a separate open obligation; $\lambda$-completeness alone is not asserted
  here to provide them.
\end{lemma}

\begin{proof}
  Choose one point in each finite solution torsor.  The differences between a
  chosen point and the restriction of the next one form the usual cocycle for
  the inverse system $(K_m)$; changing the choices changes it by a coboundary.
  Its class is $\omega$, and a coboundary correction makes the chosen points
  coherent exactly when $\omega=0$.  Under explicit surjectivity,
  Proposition~\ref{prop:synthetic-ess-coherent-tower-of-surjective}
  constructs such a tower from one initial solution.  The
  Mittag--Leffler hypothesis gives
  $\varprojlim^1K_m=0$ and hence the same conclusion.
  For finite nonempty strict fibers, use the proved finite inverse-system
  criterion directly; it does not require surjective restrictions.

  Given the stated comparisons, apply $c_R^{-1}$ to the compatible
  representative sequence.  The commuting equation diagram and the
  injectivity of $c_E$ imply that this representative has exactly the
  untruncated equation-label triple.  Conversely, the canonical
  restrictions carry an untruncated solution to a coherent tower.

  The present Lean fibers are sets of pairs in abelian groups; they have
  not been identified with homotopy fibers of mapping spaces.  Preservation
  of homotopy fibers by homotopy limits therefore does not construct the
  required comparisons for these definitions.  Deriving the displayed
  isomorphisms, their equation diagram, and their label compatibility from
  the actual $\lambda$-adic completeness equivalences remains open.
  The torsor calculation alone does not prove that model comparison or
  the untruncated extension.
\end{proof}

% SOURCE: MainPaper/main.tex:1025-1048, Proposition prop:9770ae6e.
\begin{proposition}[$\lambda^n$- and $\rho$-ESS have only $d_0$]
  \label{prop:lambda-rho-ess-only-d0}
  \uses{def:lambda-rho-delta-triangle,prop:synthetic-einfty-map-compatibility,prop:synthetic-einfty-first-quotient,prop:synthetic-ess-weightwise-transport,prop:fess-exact-middle-boundary}
  \notready
  In the extension spectral sequences of $\lambda^n$ and $\rho_{n,m}$, the
  only possibly nonzero differentials are
  $d_0^{\lambda^n}=\lambda^n$ and $d_0^\rho=\rho$.  These $d_0$ extensions have
  no crossings.  On each tridegree, $\lambda^n$ is surjective or zero and
  $\rho$ is injective or zero on $E_\infty$.
\end{proposition}

\begin{proof}
  First suppose $m<\infty$.  If one of the quotient exponents is one, use the
  mod-$\lambda$ collapse and unique-detection statement of
  Proposition~\ref{prop:synthetic-einfty-first-quotient}; when both relevant
  exponents are at least two, use the two finite $E_\infty$ formulas.  If a
  tridegree is outside the quotient range its group is zero, while in range
  $\rho$ is a cycle inclusion and $\lambda^n$ is a boundary quotient.

  If $m=\infty$, instead use the two infinite-endpoint clauses of
  Proposition~\ref{prop:synthetic-einfty-map-compatibility}: on the
  untruncated $E_\infty$ formula, $\lambda^n$ is the quotient that raises the
  boundary cutoff by $n$, and
  $\rho:\nu X\to\nu X/\lambda^n$ includes permanent representatives into the
  finite cycle cutoff (using the special-fiber clause when $n=1$).  Thus the
  sequence is exact at its middle term in every tridegree in both the finite
  and infinite branches.
  Apply the synthetic exact-middle theorem to eliminate higher differentials,
  and use the degree-zero filtration shift to obtain no-crossing.
\end{proof}

% SOURCE: MainPaper/main.tex:1068-1131, proof of Proposition prop:6de7d130.
\begin{lemma}[Shorter $\delta_n$ images are classical boundaries]
  \label{lem:delta-images-equal-classical-boundaries}
  \uses{thm:external-lambda-bockstein,prop:synthetic-einfty-map-compatibility,prop:synthetic-ess-weightwise-transport,cor:fess-triangle-shift}
  \notready
  For the infinite connecting map
  $\delta_n:\nu X/\lambda^n\to\Sigma^{1,-n}\nu X$, the subgroup generated by
  extension differentials of length less than $r$ in the target tridegree is
  exactly the image of the classical boundary subgroup
  $B_{r-1}^{s+r,t+r-1}(X)$ under the corresponding $\lambda$-multiple.
\end{lemma}

\begin{proof}
  Induct on $r$.  Write each generator of $B_{r-1}$ as the target of an
  essential classical $d_{r'}$ with $2\le r'<r$.  The Bockstein comparison and
  triangle-shift propagation turn it into a shorter $\delta_n$ extension with
  the required $\lambda$ power.  Conversely, apply the same comparison to
  every shorter $\delta_n$ extension.  The inductive hypothesis removes its
  lower-page indeterminacy.
\end{proof}

% SOURCE: MainPaper/main.tex:1068-1131, Proposition prop:6de7d130, m=infinity case.
\begin{proposition}[Infinite $\delta_n$-ESS formula]
  \label{prop:delta-ess-formula-infinite}
  \uses{lem:delta-images-equal-classical-boundaries,thm:external-lambda-bockstein,cor:fess-triangle-shift,prop:lambda-rho-ess-only-d0}
  \notready
  Suppose the classical Adams spectral sequence has
  $d_r(x)=y$, with $x\in Z_{r-1}^{s,t}$ and
  $y\in Z_\infty^{s+r,t+r-1}$.  For
  $\delta_n:\nu X/\lambda^n\to\Sigma^{1,-n}\nu X$:
  \begin{enumerate}
    \item if $r\ge n+1$, then
      $d_r^{\delta_n}(x)=\lambda^{r-n-1}y$;
    \item if $r\le n+1$, then
      $d_r^{\delta_n}(\lambda^{n+1-r}x)=y$.
  \end{enumerate}
  The formulas agree when $r=n+1$ and are equalities in the correct ESS
  quotient, whether or not the displayed extension is essential.
\end{proposition}

\begin{proof}
  Start with $n=1$, where the assertion is the $\lambda$-Bockstein theorem.
  Induct on $n$ using the commuting triangle relating multiplication by
  $\lambda$, $\delta_{n-1}$, and $\delta_n$.  Triangle propagation gives the
  formula before quotient indeterminacy.  Division by one $\lambda$ introduces
  only the next classical boundary subgroup; Lemma~\ref{lem:delta-images-equal-classical-boundaries}
  identifies and removes precisely that ESS indeterminacy.
\end{proof}

% SOURCE: MainPaper/main.tex:1068-1131, Proposition prop:6de7d130, finite m case.
% VERIFIED-IN: the printed second case places y in E_infty(nu X); the finite target requires its rho-image.
\begin{proposition}[Finite $\delta_{n,m}$-ESS formula]
  \label{prop:delta-ess-formula-finite}
  \uses{prop:delta-ess-formula-infinite,cor:fess-composition,prop:lambda-rho-ess-only-d0,prop:lambda-quotient-restriction-coherence}
  \notready
  For finite $m>n$, compose the formulas of
  Proposition~\ref{prop:delta-ess-formula-infinite} with the reduction
  $\rho:\nu X\to\nu X/\lambda^{m-n}$.  Thus the displayed target is the image
  of $\lambda^{r-n-1}y$, or of $y$ in the second case, in the finite quotient;
  it is zero once the required $\lambda$ power lies outside that quotient.
  This explicit image corrects the type ambiguity in
  \texttt{MainPaper/main.tex:1088--1094}.
\end{proposition}

\begin{proof}
  Factor $\delta_{n,m}$ as $\delta_{n,\infty}$ followed by $\rho$.  The
  $\rho$-ESS has only its injective-or-zero $d_0$, with no crossing, so ESS
  composition transports the infinite formula to the finite target.  The
  quotient formula decides exactly when the image is zero.
\end{proof}

% SOURCE: MainPaper/main.tex:1138-1144, Corollary cor:2a636737.
\begin{corollary}[$\lambda$-multiple form of the $\delta$ formula]
  \label{cor:delta-ess-lambda-multiples}
  \uses{prop:delta-ess-formula-finite}
  \notready
  In the situation above,
  \[
    d_r^\delta(\lambda^a x)=\lambda^{a+r-n-1}y
  \]
  whenever $0\le a\le n$ and
  $0\le a+r-n-1<m-n$; outside the finite quotient range the right-hand side
  and the extension are zero.
\end{corollary}

\begin{proof}
  Multiply the finite formula by $\lambda^a$, using $\lambda$-linearity of the
  synthetic Adams and extension spectral sequences.  Check the source and
  target quotient ranges directly.
\end{proof}

% SOURCE: MainPaper/main.tex:1146-1150, Remark rmk:qvoewfj.
\begin{theorem}[$\delta$-ESS and classical differential cosets are equivalent]
  \label{thm:delta-classical-equivalence}
  \uses{cor:delta-ess-lambda-multiples,lem:delta-images-equal-classical-boundaries,thm:external-lambda-bockstein}
  \notready
  For every allowed $a,n,m,r$, the equation
  $d_r^\delta(\lambda^a x)=\lambda^{a+r-n-1}y$ holds if and only if the
  classical differential coset is $d_r(x)=y$.  Both sides have the same
  indeterminacy $B_{r-1}$ after applying the stated $\lambda$ power.
\end{theorem}

\begin{proof}
  The forward direction applies the Bockstein comparison to a representative
  and uses the shorter-image lemma to identify the quotient.  The reverse
  direction is the $\lambda$-multiple formula.  Equality of the two boundary
  subgroups makes the correspondence bijective at the level of differential
  cosets, not merely at the level of nonvanishing.
\end{proof}

% SOURCE: MainPaper/main.tex:1152-1163, Example 4.13.
\begin{proposition}[$\delta_1,\delta_2,\delta_3$ regression]
  \label{prop:delta-14-stem-regression}
  \uses{prop:synthetic-14-stem-regression,cor:delta-ess-lambda-multiples}
  \notready
  The $d_2(h_4)=h_0h_3^2$ and
  $d_3(h_0h_4)=h_0d_0$ examples yield exactly the eight
  $\delta_1,\delta_2,\delta_3$ extensions and $\lambda$ powers listed in
  Example 4.13: two for $\delta_1$, three for $\delta_2$, and three for
  $\delta_3$.
\end{proposition}

\begin{proof}
  Substitute $(n,m)=(1,\infty),(2,\infty),(3,\infty)$ and each admissible
  value of $a$ into Corollary~\ref{cor:delta-ess-lambda-multiples}.  The
  synthetic $E_\infty$ regression verifies that every source and target exists
  in the displayed tridegree.
\end{proof}

% SOURCE: MainPaper/main.tex:1165-1173, Definition def:classicalcrossdiff.
\begin{definition}[Classical crossing on an $E_{n+1}$ page]
  \label{def:classical-differential-crossing}
  \uses{def:classical-adams-ss,def:ss-permanent-einfty-class}
  \notready
  Let $r\ge n+1$ and $d_r(x)=y$, with $x$ in filtration $s$.  A crossing on
  the $E_{n+1}$ page is an essential differential
  $d_{r-a-b}(x')=y'$ where the source filtration is $s+a$, the target
  filtration is $s+r-b$, $0<a\le n-1$, and
  $0\le b\le r-n-1$.
\end{definition}

% SOURCE: MainPaper/main.tex:1218-1234, Proposition prop:cross-dr-En.
\begin{proposition}[Classical crossing iff $\delta_n$ crossing]
  \label{prop:classical-crossing-iff-delta-crossing}
  \uses{def:classical-differential-crossing,def:fess-crossing,thm:delta-classical-equivalence,thm:external-synthetic-einfty-nu,thm:external-synthetic-einfty-quotient,prop:synthetic-einfty-first-quotient}
  \notready
  A classical $d_r(x)=y$ has a crossing on $E_{n+1}$ if and only if the
  corresponding extension
  $d_r^{\delta_n}(x)=\lambda^{r-n-1}y$ has a crossing in the synthetic
  $\delta_n$-ESS.
\end{proposition}

\begin{proof}
  A synthetic crossing has the unique degree form
  $d_{r-a-b}^{\delta_n}(\lambda^a x')=\lambda^{r-n-1-b}y'$.  The two
  sides are both empty when $n=1$: the classical range would require
  $0<a\le0$, while mod-$\lambda$ collapse and unique detection exclude an
  additional synthetic representative.  For $n\ge2$, the two $E_\infty$
  formulas force exactly the inequalities in
  Definition~\ref{def:classical-differential-crossing}.  Apply the
  $\delta$--classical equivalence in each direction, preserving essentiality
  through the shared boundary quotient.  The first-quotient proposition is
  used only in the $n=1$ branch.
\end{proof}

% SOURCE: MainPaper/main.tex:1236-1238, Remark nocrossE2.
\begin{corollary}[No classical crossing on $E_2$]
  \label{cor:classical-e2-no-crossing}
  \uses{def:classical-differential-crossing,prop:classical-crossing-iff-delta-crossing}
  \notready
  No classical Adams differential has a crossing on the $E_2$ page, and every
  corresponding $\delta_1$ extension has no crossing.
\end{corollary}

\begin{proof}
  Setting $n=1$ would require an integer $a$ with $0<a\le0$.  Hence no
  classical crossing exists; transfer the conclusion through
  Proposition~\ref{prop:classical-crossing-iff-delta-crossing}.
\end{proof}

% SOURCE: MainPaper/main.tex:1240-1271, Example exam:classcrossdiff.
\begin{proposition}[Stem-$38$ crossing regression]
  \label{prop:stem38-crossing-regression}
  \uses{prop:classical-crossing-iff-delta-crossing,cor:classical-e2-no-crossing,def:external-evidence}
  \notready
  Given the typed computation evidence
  $d_3(e_1)=h_1t$ and
  $d_4(h_0h_3h_5)=h_0^2x$, the former is a crossing of the latter on the
  $E_3$ and $E_4$ pages, but not on $E_2$ or on pages $E_r$ for $r\ge5$.
  The associated $\delta_n$ extensions and their weights are exactly those in
  Example 4.18.
\end{proposition}

\begin{proof}
  Check the four inequalities in the crossing definition for $n=1,2,3$ and
  $n\ge4$, then transport the two computed differential cosets using the
  crossing equivalence.  The synthetic $E_\infty$ formulas verify which
  $\lambda$ multiples remain represented.
\end{proof}

\section{Classical page-extension theorems}

% SOURCE: MainPaper/main.tex:1541-1564 and 2084-2099, the finite Ext-product
% enumeration used both for no-crossing and for page stretching.  This node is
% deliberately stated without assuming that the E_infinity extension exists.
\begin{proposition}[Hopf crossing-obstruction exclusion]
  \label{evidence:hopf-crossing-obstruction-exclusion}
  \uses{def:page-extension-crossing,def:finite-page-extension-nonliftable-crossing,thm:external-adams-one-line,def:ss-elementwise-page-presentation,def:ss-cycle-boundary-survival,def:external-evidence,evidence:normalized-hopf-detection}
  \notready
  For the candidate Hopf relation of length two from
  $h_0h_4^2$ to $h_0p$, a finite evidence record enumerates every
  degree-allowed crossing tuple and every tuple satisfying the
  nonliftable-crossing bounds of
  Definition~\ref{def:finite-page-extension-nonliftable-crossing} between
  $E_3$ and any finite page.  The list is
  empty.  A separate part enumerates every permanent source and target in the
  same forced bidegrees and excludes every actual shorter
  $(\nu_{\mathrm{Hopf}},E_\infty)$-extension; this is not inferred from the
  finite list.  Both the finite-page reductions and this permanent-candidate
  enumeration use the Hopf component of the single stored normalized
  triangle lift from Proposition~\ref{evidence:normalized-hopf-detection},
  together with its retained equality
  $\widehat{\nu_{\mathrm{Hopf}}}=[h_2]$.  Its decisive typed product check is
  \[
    h_0^2h_4^2h_2=0\ne h_0p.
  \]
  The same record supplies the independent typing facts
  \[
    h_0h_4^2\in Z_\infty(S^3),
    \qquad h_0p\in Z_\infty(S^0).
  \]
  The first follows from the known $d_2(h_5)=h_0h_4^2$ and
  $B_2\subseteq Z_\infty$; the second is a complete outgoing-differential
  exclusion in the target bidegree.  Thus both the source and target remain
  typed on every later page.  It also certifies
  $h_0p\notin B_2$ and that no shorter Hopf-ESS image reaches its complete
  target coset.  Consequently, whenever the candidate length-two relation
  exists at a finite or infinite level, it is essential.
  Thus the record contains both a finite-page obstruction list and an
  independently computed untruncated candidate list.  Neither branch assumes
  nor concludes that the candidate relation already exists at its respective
  finite or infinite level.
\end{proposition}

\begin{proof}
  Enumerate the finitely many integers $a,a',b$ allowed by the crossing and
  nonliftable-crossing inequalities and then enumerate the Ext groups in their
  forced source and target bidegrees.  Intersect the same finite bidegree lists
  with $Z_\infty$ to obtain the independent untruncated-candidate list.  The
  stored equality $\widehat{\nu_{\mathrm{Hopf}}}=[h_2]$ identifies, on every
  finite reduction and on the untruncated list, the only possible shorter Hopf
  product with the displayed zero class, whereas the proposed target is
  nonzero.  Store
  the degree bounds, cycle-range checks, complete incoming/outgoing target
  search, and typed product inequality in the evidence record.  The source cycle fact is
  also obtained from the one-line differential and the inclusion of boundaries
  into permanent representatives.  No inverse-limit representative or
  $E_\infty$ extension is used in this enumeration.
\end{proof}

% SOURCE: MainPaper/main.tex:1284-1298, E_0 formulas before Definition def:6c076a33.
\begin{proposition}[$E_0$ presentation for the reduced-map ESS]
  \label{prop:hat-map-ess-e0}
  \uses{def:hat-map-page-reduction,def:synthetic-extension-ss,thm:external-synthetic-einfty-nu,thm:external-synthetic-einfty-quotient,prop:synthetic-einfty-first-quotient}
  \notready
  For finite $r$, in weight $t-k+e(f)$, the $E_0$ page of the
  $\widehat f_{r-1}$-ESS is
  \[
  \frac{Z_{r-1-k}^{s,t}(X)}{B_{1+k}^{s,t}(X)}\oplus
  \frac{Z_{r-1-k+e(f)}^{s,t}(Y)}{B_{1+k-e(f)}^{s,t}(Y)}.
  \]
  This displayed formula is for finite $r$.  At $r=\infty$, the $E_0$ page
  of the $\widehat f$-ESS is instead
  \[
  \frac{Z_\infty^{s,t}(X)}{B_{1+k}^{s,t}(X)}\oplus
  \frac{Z_\infty^{s,t}(Y)}{B_{1+k-e(f)}^{s,t}(Y)};
  \]
  no expression such as $\infty-1-k$ is formed.
  For finite $r$, its length-$n$ differential is a map from
  \[
    {}^{\widehat f_{r-1}}\!Z_{n-1}^{s,t,t-k}(X)
      \subseteq Z_{r-1-k}^{s,t}(X)/B_{1+k}^{s,t}(X)
  \]
  to
  \[
    \frac{Z_{r-1-k-n+e(f)}^{s+n,t+n}(Y)/
      B_{1+k+n-e(f)}^{s+n,t+n}(Y)}
    {{}^{\widehat f_{r-1}}\!B_{n-1}^{s+n,t+n,t-k}(Y)}.
  \]
  It vanishes for degree reasons when $n<\operatorname{AF}(f)$ or
  $n>r-2-k+e(f)$.
  At $r=\infty$, replace $\widehat f_{r-1}$ by $\widehat f$ and every cycle
  numerator in these source and target groups by $Z_\infty$; only the lower
  vanishing bound $n<\operatorname{AF}(f)$ remains.
\end{proposition}

\begin{proof}
  If $r=2$, substitute the source and target weights into the first-quotient
  edge presentation.  If $3\le r<\infty$, use the finite-quotient
  $E_\infty$ formula at exponent $r-1\ge2$.  If $r=\infty$, use the
  untruncated $Z_\infty/B_j$ formula separately.  Then apply the elementwise ESS presentation,
  simplify the two cycle and boundary subscripts over integers, and check the
  quotient range before converting to a finite page index.
\end{proof}

% SOURCE: MainPaper/main.tex:1391-1402, Example exam:extonEn(1).
\begin{proposition}[Hopf-map $E_2$ page extension]
  \label{prop:hopf-e2-page-extension}
  \uses{def:page-extension,evidence:normalized-hopf-detection,prop:synthetic-einfty-first-quotient,def:external-evidence}
  \notready
  Let $\nu_{\mathrm{Hopf}}:S^3\to S^0$ denote the classical Hopf map.  Typed Ext
  product evidence supplies the essential extension
  \[
    d_1^{\nu_{\mathrm{Hopf}},E_2}(h_5)=h_2h_5.
  \]
  The notation $\nu_{\mathrm{Hopf}}$ is kept distinct from the synthetic functor
  $\nu$.
\end{proposition}

\begin{proof}
  Reduce the precise lifted-triangle component identified in
  Proposition~\ref{evidence:normalized-hopf-detection} modulo $\lambda$.
  The first-quotient uniqueness statement identifies its reduction with
  multiplication by $h_2$, so the stated extension is exactly the supplied
  Ext product.  Thus this node and the later $E_3/E_\infty$ nodes use the same
  chosen normalized Hopf map.
\end{proof}

% SOURCE: MainPaper/main.tex:1403-1423, Example exam:extonEn(2); proved later by Example exam:Mahowald.
\begin{proposition}[Hopf-map $E_3$ page extension]
  \label{prop:hopf-e3-page-extension}
  \uses{def:page-extension,thm:generalized-mahowald,prop:mahowald-cofiber-regression,evidence:hopf-crossing-obstruction-exclusion}
  \notready
  The normalized Hopf map has the essential extension
  \[
    d_2^{\nu_{\mathrm{Hopf}},E_3}(h_0h_4^2)=h_0p.
  \]
  Equivalently, modulo $\lambda^2$ there are detected representatives with
  $[h_0h_4^2][h_2]=[\lambda h_0p]$.
\end{proposition}

\begin{proof}
  This is the conclusion of the cofiber regression obtained from the
  generalized Mahowald trick.  Transport that synthetic extension through
  Definition~\ref{def:page-extension} and verify the displayed bidegrees.  The
  independent Hopf obstruction certificate says that $h_0p$ is nonzero modulo
  $B_2$ and is not a shorter Hopf-ESS image, so the complete target coset does
  not contain zero; hence the extension is essential.
\end{proof}

% SOURCE: MainPaper/main.tex:1424-1433, Example exam:extonEn(3); proved later by Example exam:stretchext.
\begin{proposition}[Hopf-map $E_\infty$ page extension]
  \label{prop:hopf-einfty-page-extension}
  \uses{prop:hopf-e3-page-extension,prop:page-stretch-regression}
  \notready
  The preceding relation lifts to an essential extension
  \[
    d_2^{\nu_{\mathrm{Hopf}},E_\infty}(h_0h_4^2)=h_0p.
  \]
\end{proposition}

\begin{proof}
  This is the intended conclusion of the page-stretch regression.
  Its untruncated lifting argument remains to be completed; after existence
  is proved, the independent nonzero target-coset clause gives essentiality.
\end{proof}

% SOURCE: MainPaper/main.tex:1391-1433, Example exam:extonEn(4).
\begin{proposition}[Multiplication-by-two page-extension regression]
  \label{prop:two-page-extension-regression}
  \uses{def:page-extension-coset,prop:no-crossing-iff-uniform-detection,def:external-evidence}
  \notready
  Given the stated $14$-stem evidence and the comparison
  $\lambda[h_0]=2$, multiplication by two has
  $d_2^{2,E_\infty}(h_0h_3^2)=0$, while
  $d_1^{2,E_\infty}(d_0)=h_0d_0$ is essential and makes the alternative target
  $d_2^{2,E_\infty}(h_0h_3^2)=h_0d_0$ inessential.
\end{proposition}

\begin{proof}
  Evaluate multiplication by $[h_0]$ on the supplied synthetic homotopy group.
  Its only nonzero leading target comes from the higher-filtration class
  $[\lambda d_0]$.  The representative and inessentiality criteria yield the
  three asserted page-extension statuses.
\end{proof}

% SOURCE: MainPaper/main.tex:1515-1539, Proposition prop:cross-f-Er.
\begin{proposition}[Page crossing iff synthetic crossing]
  \label{prop:page-crossing-iff-synthetic-crossing}
  \uses{def:page-extension-crossing,def:fess-crossing,prop:synthetic-einfty-map-compatibility,prop:synthetic-ess-weightwise-transport}
  \notready
  An $(f,E_r)$-extension has a page crossing if and only if its defining
  $\widehat f_{r-1}$-ESS extension has a crossing for finite $r$.  For
  $r=\infty$, the equivalent objects are the actual untruncated
  $(f,E_\infty)$ crossing and an actual crossing in the $\widehat f$-ESS.
\end{proposition}

\begin{proof}
  A synthetic crossing necessarily has the form
  $d_{n-a-b}^{\widehat f_{r-1}}(\lambda^ax')
   =\lambda^{n-b-e(f)}y'$.  Compare the quotient of exponent $r-1$ with that
  of exponent $r-1-a$.  Dividing by $\lambda^a$ introduces exactly
  $B_{1+n-b-e(f)}$; therefore essentiality is equivalent to the nonboundary
  condition in the finite page-crossing definition.

  At $r=\infty$, start with an actual untruncated shorter page extension on
  permanent representatives.  The infinite-endpoint clause of
  Proposition~\ref{prop:synthetic-einfty-map-compatibility} identifies
  multiplication by $\lambda^a$ with the corresponding boundary-cutoff
  quotient on $Z_\infty/B_j$.  This converts that actual page witness to an
  actual crossing in the untruncated $\widehat f$-ESS, and the same quotient
  calculation gives the converse.  No finite solution tower is promoted to an
  infinite witness in this argument.
\end{proof}

% SOURCE: MainPaper/main.tex:1541-1579, Example exam:Ercross.
\begin{proposition}[Page-crossing regression]
  \label{prop:page-crossing-regression}
  \uses{prop:hopf-e2-page-extension,prop:hopf-e3-page-extension,prop:hopf-einfty-page-extension,evidence:hopf-crossing-obstruction-exclusion,prop:two-page-extension-regression,prop:page-crossing-iff-synthetic-crossing,def:external-evidence}
  \notready
  The four examples of Section 5 have the stated crossing behavior: the
  $E_2$, $E_3$, and $E_\infty$ Hopf-map extensions have no crossings, while
  the zero multiplication-by-two extension is crossed by the essential
  length-one $d_0$ target.  The no-crossing certificates include the Ext
  relation $h_0^2h_4^2h_2=0\ne h_0p$.
\end{proposition}

\begin{proof}
  There is no positive $a$ in the $E_2$ crossing range.  For the actual
  Hopf-map $E_3$ and $E_\infty$ extensions, apply the independent finite
  obstruction certificate.  Its separate permanent-candidate enumeration
  excludes every actual untruncated shorter extension without having assumed
  the candidate $E_\infty$ relation.  For multiplication by two, the essential
  $d_1$ extension from $d_0$ meets the crossing inequalities and witnesses the
  failure of no-crossing.
\end{proof}

% SOURCE: MainPaper/main.tex:1581-1597, unlabelled Proposition 5.14.
\begin{proposition}[$E_\infty$ page extension gives a classical extension]
  \label{prop:einfty-page-to-classical-extension}
  \uses{def:page-extension,thm:external-lambda-inversion,def:reindexed-ss-morphism,def:f-extension-essential}
  \notready
  If $x\in Z_\infty^{s,t}(X)$ and
  $y\in Z_\infty^{s+n,t+n}(Y)$ satisfy
  $d_n^{f,E_\infty}(x)=y$, then after passing to classical $E_\infty$ cosets,
  \[
    d_n^f(x+B_\infty)=y+B_\infty.
  \]
  The resulting classical extension may be inessential.
\end{proposition}

\begin{proof}
  Invert $\lambda$ in the $\widehat f$-ESS.  Generic-fiber equivalence
  identifies its source and target with the classical $f$-ESS and sends the
  normalized map to $f$.  Naturality of the induced reindexed page maps sends
  the synthetic differential coset to the displayed classical coset.
\end{proof}

\section{Generalized Leibniz and Mahowald rules}

% SOURCE: MainPaper/main.tex:1606-1652, Theorem thm:e73f481e.
\begin{theorem}[Generalized Leibniz rule]
  \label{thm:generalized-leibniz}
  \uses{prop:delta-ess-formula-infinite,thm:delta-classical-equivalence,prop:classical-crossing-iff-delta-crossing,def:page-extension,prop:page-crossing-iff-synthetic-crossing,cor:fess-square-naturality}
  \lean{KIP126.Kervaire.Route.Tools.GeneralizedLeibnizLaw}
  \notready
  Let $f:X\to Y$, $2\le n\le r$,
  $e(f)\le m\le n-2+e(f)$, and $l\ge e(f)$.  Suppose
  \[
  \begin{aligned}
    &x\in Z_{r-1}^{s,t}(X),&&
    y\in Z_{r-1-m+e(f)}^{s+m,t+m}(Y),\\
    &x_\infty\in Z_\infty^{s+r,t+r-1}(X),&&
    y_\infty\in Z_\infty^{s+r+l,t+r+l-1}(Y),
  \end{aligned}
  \]
  and assume:
  \begin{enumerate}
    \item $d_r(x)=x_\infty$;
    \item $d_m^{f,E_n}(x)=y$;
    \item $d_l^{f,E_\infty}(x_\infty)=y_\infty$;
    \item either the first differential has no crossing on $E_n$, or the
      second page extension has no crossing;
    \item the third page extension has no crossing.
  \end{enumerate}
  Then $d_{r+l-m}(y)=y_\infty$.
  The linked Lean declaration specifies this law for a fixed actual
  normalized page family; its validity is still a delivery obligation.
  It does not assert that arbitrary comparison isomorphisms satisfy the law.
\end{theorem}

\begin{proof}
  Form the square whose horizontal maps are
  $\widehat f_{n-1}$ and $\widehat f$ and whose vertical maps are the two
  connecting maps from reduction modulo $\lambda^{n-1}$.  Translate (1) by the
  $\delta$ formula and (2)--(3) by the definition of page extension.  The two
  crossing comparison theorems turn hypothesis (4) into synchronized-source
  no-crossing and hypothesis (5) into target no-crossing.  Apply ESS square
  propagation.  Finally use the $\delta$--classical equivalence to remove the
  displayed $\lambda$ powers and obtain the asserted classical differential
  coset.
\end{proof}

% SOURCE: MainPaper/main.tex:1664-1692, Example exam:Leibnizyes.
\begin{proposition}[Positive Leibniz regression]
  \label{prop:leibniz-positive-regression}
  \uses{thm:generalized-leibniz,thm:external-adams-one-line,prop:hopf-e2-page-extension,prop:hopf-einfty-page-extension,cor:classical-e2-no-crossing,prop:page-crossing-regression}
  \notready
  The inputs
  $d_2(h_5)=h_0h_4^2$,
  $d_1^{\nu_{\mathrm{Hopf}},E_2}(h_5)=h_2h_5$, and
  $d_2^{\nu_{\mathrm{Hopf}},E_\infty}(h_0h_4^2)=h_0p$
  satisfy the generalized Leibniz hypotheses and give
  \[
    d_3(h_2h_5)=h_0p.
  \]
\end{proposition}

\begin{proof}
  Substitute $(n,r,m,l,e(f))=(2,2,1,2,1)$.  The $E_2$ classical crossing
  condition is automatic, and the two Hopf-map page extensions have the
  no-crossing certificates from the regression.  The theorem gives a
  differential of length $2+2-1=3$ with the claimed target.
\end{proof}

% SOURCE: MainPaper/main.tex:1694-1720, Example exam:Leibnizno.
\begin{proposition}[No-crossing is necessary]
  \label{prop:leibniz-negative-regression}
  \uses{thm:generalized-leibniz,prop:two-page-extension-regression,prop:page-crossing-regression,def:external-evidence}
  \notready
  For $f=2$ and the computed differential
  $d_2(h_4)=h_0h_3^2$, all generalized Leibniz hypotheses except target
  no-crossing can be met using
  $d_2^{2,E_\infty}(h_0h_3^2)=0$.  The missing hypothesis is false because
  $d_1^{2,E_\infty}(d_0)=h_0d_0$ crosses it.  Omitting that hypothesis would
  assert the false zero differential $d_3(h_0h_4)=0$.
\end{proposition}

\begin{proof}
  Evaluate every hypothesis using the page-extension regression and the typed
  classical computation.  The length-one extension meets the crossing bounds,
  so the generalized theorem is inapplicable.  The known nonzero differential
  on $h_0h_4$ supplies the regression check that the weakened statement is
  false.
\end{proof}

% SOURCE: MainPaper/main.tex:1733-1753, Remark rem:chuaerror; Chua Definition 12.5 and Theorem 12.9.
\begin{proposition}[Counterexample to the maximal-extension rule]
  \label{prop:chua-rule-counterexample}
  \uses{prop:leibniz-negative-regression,def:external-result,def:external-evidence}
  \notready
  The maximal-$[h_0]$ lift of $h_0h_3^2$ in the stated synthetic $14$-stem
  example maps to zero, while another lift maps to $[h_0d_0]$.  Consequently
  the rule quoted as Chua Theorem 12.9, when read without a no-crossing
  hypothesis, predicts the false conclusion of
  Proposition~\ref{prop:leibniz-negative-regression}.
\end{proposition}

\begin{proof}
  Use the two distinct detected lifts and their products from the supplied
  evidence.  The zero product is maximally $\lambda$-divisible under the cited
  definition.  Substitution into the quoted rule yields the false zero target,
  so these data form a concrete counterexample; no claim about a corrected
  version is made.
\end{proof}

% SOURCE: MainPaper/main.tex:1780-1917, complete hypotheses and degree package for Theorem thm:158d451a.
\begin{definition}[Generalized Mahowald input package]
  \label{def:mahowald-input-data}
  \uses{prop:normalized-synthetic-triangle,def:page-extension,def:classical-differential-crossing}
  \notready
  A Mahowald input package consists of a distinguished triangle
  $X\xrightarrow fY\xrightarrow gZ\xrightarrow h\Sigma X$ with
  $e(f)+e(g)+e(h)=1$, integers $s,t,n,m,l$, and
  \[
    r=n+m+l,\quad n_1=n-e(f)\ge1,\quad
    m_1=m-e(g)\ge0,\quad l_1=l-e(h)\ge0,\quad
    r'=r-m_1=n_1+l_1+1.
  \]
  It contains classes
  \[
  \begin{aligned}
    x&\in Z_{n_1}^{s+l,t+l-1}(X),&
    y&\in Z_{m_1+1}^{s+n+l,t+n+l-1}(Y),\\
    \bar x&\in Z_{r-1}^{s,t}(Z),&
    \bar y&\in Z_\infty^{s+r,t+r-1}(Z),
  \end{aligned}
  \]
  proofs of $d_l^{h,E_{r'}}(\bar x)=x$, $d_r(\bar x)=\bar y$,
  $d_m^{g,E_{m_1+2}}(y)=\bar y$, and a proof that either the first page
  extension has no crossing or the classical differential has no crossing on
  $E_{r'}$.
  In the Lean statement these are arguments of
  \texttt{KIP126.Kervaire.Route.Tools.GeneralizedMahowaldLaw}, not freely chosen page
  operations. The target of the $h$-extension is an actual $E_2(\Sigma X)$
  label; the following comparison relates it to the displayed $X$ label.
\end{definition}

\begin{definition}[Actual tower suspension comparison]
  \label{def:mahowald-tower-suspension-comparison}
  \uses{def:hf2-adams-tower,def:ss-elementwise-page-presentation}
  \lean{KIP126.Classical.Adams.Suspension.TowerComparison}
  \lean{KIP126.Classical.Adams.Suspension.TowerComparison.DesuspendsClass}
  \notready
  A suspension comparison identifies every term of the actual Adams tower
  of $\Sigma X$ with the suspension of the corresponding term for $X$,
  is the identity at stage zero, and commutes with all tower maps. Its
  layer isomorphisms commute with the actual cofiber inclusions and with
  the connecting maps, including the suspension sign. Two $E_2$ labels
  are related when they have raw second-cycle representatives whose
  actual first-page maps agree after this desuspension and whose quotient
  classes are the specified labels. Existence and independence of this
  comparison have not been constructed; an arbitrary $E_2$ equivalence
  is not a substitute for these diagrams.
\end{definition}

% SOURCE: MainPaper/main.tex:1780-1917, representative choices in the proof of Theorem thm:158d451a.
\begin{lemma}[Synchronized Mahowald representatives]
  \label{lem:mahowald-synchronized-representatives}
  \uses{def:mahowald-input-data,prop:no-crossing-iff-uniform-detection,prop:classical-crossing-iff-delta-crossing,prop:page-crossing-iff-synthetic-crossing,thm:delta-classical-equivalence}
  \notready
  For a Mahowald input package, one can choose
  \[
    [\bar x]\in\pi_{t-s,t}(\nu Z/\lambda^{r'-1})
  \]
  detected by $\bar x$ such that
  $\widehat h_{r'-1}[\bar x]$ is detected by
  $\lambda^{l_1}x$ in
  $\pi_{t-s-1,t-1+e(h)}(\nu X/\lambda^{r'-1})$ and
  $\delta_{r'-1}[\bar x]$ is detected by
  $\lambda^{m_1}\bar y$ in
  $\pi_{t-s-1,t+r'-1}(\nu Z/\lambda^{m_1+1})$.
\end{lemma}

\begin{proof}
  Translate the page extension and classical differential into the two
  synthetic ESS extensions in the common quotient.  The assumed no-crossing
  of at least one extension invokes uniform detection for every representative
  chosen to realize the other.  Hence one representative realizes both
  leading terms simultaneously.
\end{proof}

% SOURCE: MainPaper/main.tex:1886-1907, the representative for y, degree-uniqueness, and the application of May's pullback lemma.
\begin{lemma}[May-compatible lift in the Mahowald diagram]
  \label{lem:mahowald-compatible-quotient-lift}
  \uses{def:mahowald-input-data,lem:mahowald-synchronized-representatives,thm:external-may-smash-boundary,prop:synthetic-ess-weightwise-transport,prop:fextension-iff-detected,thm:external-synthetic-einfty-quotient,prop:synthetic-einfty-first-quotient,prop:normalized-synthetic-triangle,def:lambda-rho-delta-triangle}
  \notready
  For a Mahowald input package, choose the synchronized representative
  $[\bar x]$ of Lemma~\ref{lem:mahowald-synchronized-representatives}.
  There are a representative

  \[
    [y]\in\pi_{t-s-1,t+n+l-1}(\nu Y/\lambda^{m_1+1})
  \]
  detected by $y$ and a class

  \[
    u\in\pi_{t-s-1,t-1+e(h)}(\nu X/\lambda^r)
  \]
  detected by $\lambda^{l_1}x$ such that

  \[
    \rho(u)=\widehat h_{r'-1}([\bar x]),
    \qquad
    \widehat f_r(u)=\lambda^{r'-1}[y].
  \]
\end{lemma}

\begin{proof}
  The fourth page-extension hypothesis and the representative criterion choose
  $[y]$ whose image under $\widehat g_{m_1+1}$ is detected by
  $\lambda^{m_1}\bar y$.  If $m_1=0$, the first-quotient collapse and its
  unique-detection conclusion make that mod-$\lambda$ homotopy class unique;
  if $m_1\ge1$, the finite-quotient $E_\infty$ formula gives uniqueness in the
  indicated top weight.  Hence in both cases
  $\widehat g_{m_1+1}([y])$ equals
  $\delta_{r'-1}([\bar x])$.  Apply May's smash-boundary theorem to the
  normalized synthetic triangle and the $\lambda$--$\rho$--$\delta$ triangle.
  Its pullback class is $u$, and its two structure maps give exactly the two
  displayed equalities.
\end{proof}

% SOURCE: MainPaper/main.tex:1900-1910, equation eq:f4848f29 and the ensuing one-way survival deduction.
\begin{lemma}[A compatible quotient lift gives source survival]
  \label{lem:mahowald-quotient-lift-survival}
  \uses{def:mahowald-input-data,lem:mahowald-synchronized-representatives,thm:external-synthetic-einfty-quotient,prop:synthetic-einfty-map-compatibility}
  \notready
  For a Mahowald input package, suppose

  \[
    u\in\pi_{t-s-1,t-1+e(h)}(\nu X/\lambda^r)
  \]
  is detected by $\lambda^{l_1}x$ and its $\rho$-image is the synchronized
  representative detected by $\lambda^{l_1}x$ modulo $\lambda^{r'-1}$.
  Then

  \[
    x\in Z_{n+m+e(h)}^{s+l,t+l-1}(X).
  \]
  No converse for a preselected representative in the smaller quotient is
  asserted.
\end{lemma}

\begin{proof}
  Apply the finite-quotient $E_\infty$ formula to the class $u$.  Its exponent
  after removing the displayed $\lambda^{l_1}$ is
  $r-l_1=n+m+e(h)$, so the existence of $u$ places $x$ in the corresponding
  cycle subgroup.  Compatibility of $\rho$ verifies that its reduction is the
  synchronized smaller-quotient representative; it is not used to manufacture
  a lift from survival alone.
\end{proof}

% SOURCE: MainPaper/main.tex:1903-1917, final lambda-division step of Theorem thm:158d451a.
\begin{lemma}[$\lambda$-division indeterminacy in the Mahowald diagram]
  \label{lem:mahowald-lambda-division}
  \uses{def:mahowald-input-data,prop:synthetic-einfty-map-compatibility,def:page-extension-coset}
  \notready
  For a Mahowald input package, if the smash diagram gives
  \[
    d_n^{\widehat f_r}(\lambda^{l_1}x)=\lambda^{r'-1}y,
  \]
  then division by $\lambda^{l_1}$ is defined in the quotient of exponent
  $r-l_1=n+m+e(h)$ and changes the target exactly by
  $B_{r'}^{s+n+l,t+n+l-1}(Y)$.  The result is
  \[
    d_n^{f,E_{n+m+1+e(h)}}(x)
       \equiv y\pmod{B_{r'}^{s+n+l,t+n+l-1}(Y)}.
  \]
\end{lemma}

\begin{proof}
  Apply the quotient-map description of multiplication by $\lambda^{l_1}$.
  Its kernel is the boundary cutoff enlarged by $l_1$; after substituting the
  definitions of $n_1,m_1,l_1,r'$, this cutoff is $B_{r'}$.  Compare with the
  complete page-extension target coset to obtain the asserted equivalence.
\end{proof}

% SOURCE: MainPaper/main.tex:1780-1917, Theorem thm:158d451a.
\begin{theorem}[Generalized Mahowald trick]
  \label{thm:generalized-mahowald}
  \uses{def:mahowald-input-data,def:mahowald-tower-suspension-comparison,prop:normalized-synthetic-triangle,def:lambda-rho-delta-triangle,lem:mahowald-synchronized-representatives,lem:mahowald-compatible-quotient-lift,lem:mahowald-quotient-lift-survival,lem:mahowald-lambda-division}
  \lean{KIP126.Kervaire.Route.Tools.GeneralizedMahowaldLaw}
  \notready
  Let
  $X\xrightarrow fY\xrightarrow gZ\xrightarrow h\Sigma X$ be distinguished
  with $e(f)+e(g)+e(h)=1$.  Put
  $r=n+m+l$, $n_1=n-e(f)\ge1$,
  $m_1=m-e(g)\ge0$, $l_1=l-e(h)\ge0$, and
  $r'=r-m_1=n_1+l_1+1$.  Suppose
  \[
  \begin{aligned}
    x&\in Z_{n_1}^{s+l,t+l-1}(X),&
    y&\in Z_{m_1+1}^{s+n+l,t+n+l-1}(Y),\\
    \bar x&\in Z_{r-1}^{s,t}(Z),&
    \bar y&\in Z_\infty^{s+r,t+r-1}(Z),
  \end{aligned}
  \]
  and assume
  \begin{enumerate}
    \item $d_l^{h,E_{r'}}(\bar x)=x$;
    \item $d_r(\bar x)=\bar y$;
    \item either (1) has no page crossing or (2) has no crossing on $E_{r'}$;
    \item $d_m^{g,E_{m_1+2}}(y)=\bar y$.
  \end{enumerate}
  Then $x\in Z_{n+m+e(h)}$ and
  \[
    d_n^{f,E_{n+m+1+e(h)}}(x)
       \equiv y\quad\text{modulo }B_{r'}.
  \]
  The linked Lean declaration is the parameterized law, using one actual
  distinguished triangle, its three normalized page families, and the
  tower suspension comparison above. A proof of this law remains to be
  supplied; defining the predicate is not a proof of the theorem.
\end{theorem}

\begin{proof}
  Smash the normalized synthetic triangle with the quotient triangle of
  exponents $r$ and $r'-1$.  The synchronization lemma chooses the common
  representative of $\bar x$.  The May-compatible-lift lemma uses condition
  (4), the degree-unique common target, and May's pullback theorem to produce a
  single class mapping to both sides of the diagram.  The one-way
  quotient-lift lemma turns that actual $\rho$ lift into the stronger survival
  statement for $x$.  Its other image is the normalized $f$ relation; divide
  by $\lambda^{l_1}$ using the final lemma to obtain the page extension with
  exactly $B_{r'}$ indeterminacy.
\end{proof}

% SOURCE: MainPaper/main.tex:1930-1977, Example exam:Mahowald; computed cofiber d_2 is external evidence.
\begin{proposition}[Cofiber regression for the Mahowald trick]
  \label{prop:mahowald-cofiber-regression}
  \uses{thm:generalized-mahowald,lem:normalized-hopf-exactness-case,def:hopf-cofiber-cell-notation,cor:classical-e2-no-crossing,def:external-evidence}
  \notready
  Given the computed differential
  $d_2(h_0h_4^2[4])=h_0p[0]$ in the Adams spectral sequence of
  $S^0/\nu_{\mathrm{Hopf}}$ and the two displayed Ext edge-map facts, the
  generalized Mahowald trick proves that $h_0h_4^2$ survives to $E_3$ and
  \[
    d_2^{\nu_{\mathrm{Hopf}},E_3}(h_0h_4^2)=h_0p.
  \]
\end{proposition}

\begin{proof}
  Use Lemma~\ref{lem:normalized-hopf-exactness-case} to instantiate the Hopf
  triangle with $(e(f),e(g),e(h))=(1,0,0)$ and its one chosen lifted triangle.
  Substitute $m=l=0$, $n=r=r'=2$, $n_1=1$, and $m_1=l_1=0$ into the theorem.
  Bottom- and top-cell Ext exactness give the two length-zero page extensions;
  the supplied cofiber differential gives condition (2), and $E_2$
  no-crossing gives condition (3).  The theorem produces the stated
  extension with zero additional indeterminacy in this degree.
\end{proof}

% SOURCE: MainPaper/main.tex:1981-2068, Proposition prop:dec738d3, first part.
\begin{proposition}[Restriction of a page extension]
  \label{prop:page-extension-restrict}
  \uses{def:page-extension,cor:fess-map-of-stable-pages,prop:lambda-rho-ess-only-d0,prop:lambda-quotient-restriction-coherence}
  \notready
  If $d_n^{f,E_r}(x)=y$ and $2\le r'<r$, reduction along
  \[
    \rho_X:\nu X/\lambda^{r-1}\to\nu X/\lambda^{r'-1},
    \qquad
    \rho_Y:\nu Y/\lambda^{r-1}\to\nu Y/\lambda^{r'-1}
  \]
  gives
  $d_n^{f,E_{r'}}(x)=y$ whenever the source and target are defined.
\end{proposition}

\begin{proof}
  The square of reduced normalized maps commutes.  Since both vertical
  $\rho$-ESS have only $d_0$, the stable-page map theorem induces a morphism of
  the two $\widehat f$ extension spectral sequences.  Naturality sends the
  length-$n$ differential to its earlier-page version.
\end{proof}

\begin{lemma}[Cokernel criterion for a specified strict solution lift]
  \label{lem:strict-solution-cokernel-obstruction}
  \uses{def:strict-ess-representative-fiber,def:filtered-chain-map}
  \lean{KIP126.Core.SpectralSequence.FilteredComplex.Solutions.liftingMap,KIP126.Core.SpectralSequence.FilteredComplex.Solutions.LiftingObstruction,KIP126.Core.SpectralSequence.FilteredComplex.Solutions.obstruction,KIP126.Core.SpectralSequence.FilteredComplex.Solutions.obstruction_eq_zero_iff}
  \leanok
  Let $f:C\to D$ be a filtered chain map, $R_C,R_D$ the actual groups
  of representative pairs, and $E_C$ the equation group of
  Definition~\ref{def:strict-ess-representative-fiber}.
  With $\rho_R:R_C\to R_D$ the induced representative map, set
  \[
    L(z)=(\Phi_C(z),\rho_R(z)),\qquad
    \operatorname{Ob}(f)=\operatorname{coker}\bigl(L:R_C\to E_C\times R_D\bigr).
  \]
  For a specified earlier solution $a$ with the images of labels $(x,y)$,
  its obstruction is the class of $((x,y,0),a)$.  That class is zero
  exactly when there is a strict solution $b$ with labels $(x,y)$ and
  $\rho_R(b)=a$.  No nonemptiness of the later fiber is assumed.
  This criterion concerns that particular representative pair $a$.
  The paper's later-relation conclusion has a different quantifier; this
  cokernel criterion does not imply a no-crossing criterion.
\end{lemma}
\begin{proof}
  \leanok
  Zero in the cokernel means that $((x,y,0),a)$ lies in the image of $L$.
  Its two coordinates are exactly the later equation and the prescribed
  restriction.  Conversely, any such solution supplies the required preimage.
\end{proof}

\begin{lemma}[Surjectivity on strict fibers and their difference groups]
  \label{lem:strict-solution-restriction-surjectivity}
  \uses{def:strict-ess-representative-fiber,def:filtered-chain-map,lem:strict-solution-cokernel-obstruction}
  \lean{KIP126.Core.SpectralSequence.FilteredComplex.Solutions.restrict_displacement,KIP126.Core.SpectralSequence.FilteredComplex.Solutions.restrict_surjective_iff_differences_surjective,KIP126.Core.SpectralSequence.FilteredComplex.Solutions.forall_obstruction_eq_zero_iff_surjective,KIP126.Core.SpectralSequence.FilteredComplex.Solutions.forall_obstruction_eq_zero_iff_differences_surjective}
  \leanok
  Let $f:C\to D$ be a filtered chain map and fix source and target
  labels for its actual strict solution fibers.  Restriction is surjective
  on these fibers if and only if the specified lifting obstruction vanishes
  for every earlier solution.  This equivalence needs no nonemptiness
  assumption.

  If a later solution $a$ is specified, surjectivity on the strict fibers is
  also equivalent to surjectivity of the actual induced homomorphism
  \[
    K_C=\ker\Phi_C\longrightarrow K_D=\ker\Phi_D.
  \]
  In particular, existence of one later solution does not establish either
  surjectivity claim.  The theorem reduces the remaining lifting obligation
  to a concrete homomorphism on the homogeneous solution groups.
\end{lemma}
\begin{proof}
  \leanok
  The first equivalence is the cokernel zero criterion applied to every
  earlier solution.  Relative to $a$ and its restriction, the coordinate
  equivalences identify a solution with its displacement from the chosen
  basepoint.  Restriction preserves displacements and translations, so
  lifting all solutions is equivalent to lifting all displacements.
\end{proof}

% SOURCE: MainPaper/main.tex:1437-1441,
% Definition def:41d51149; 1515-1536, Proposition prop:cross-f-Er; and
% 1981-2081, Proposition prop:dec738d3 and Corollary cor:dfc6043e.
% This definition includes b=0 and uses the larger ordinary boundary cutoff.
\begin{definition}[Finite nonliftable crossing candidate]
  \label{def:finite-page-extension-nonliftable-crossing}
  \uses{def:page-extension,def:page-extension-crossing,def:page-extension-coset}
  \lean{KIP126.Kervaire.Route.FiniteExtensionNonliftableCrossing}
  \leanok
  Fix a normalized family $P$, finite pages $r',r$, and integers $n,s,t$.
  A nonliftable crossing candidate consists of natural numbers $a,b$ with
  \[
    0<a\le r'-2,\qquad 0\le b\le n-a-e(f),
  \]
  classes
  \[
    x'\in Z_{r'-1-a}^{s+a,t+a}(X)
       \setminus Z_{r-1-a}^{s+a,t+a}(X),\qquad
    y'\in Z_{r'-1-n+b+e(f)}^{s+n-b,t+n-b}(Y),
  \]
  and an actual essential shorter extension in the same family,
  \[
    d_{n-a-b}^{f,E_{r'-a}}(x')=y',\qquad
    y'\notin B_{1+n-b-e(f)}^{s+n-b,t+n-b}(Y).
  \]
  The last exclusion is the larger ordinary boundary cutoff from the
  crossing definition; essentiality alone does not supply it.
  The source and target cycle conditions are those of the actual shorter
  extension witness.  The bound on $b$ also gives $a\le n-e(f)$.

  This predicate does not relate $y'$ to a chosen original target $y$.
  It enumerates all candidates in these degrees and is used only through
  their exclusion.  It is not an equivalence characterizing loss of
  essentiality, and does not describe the cokernel of a raw strict-solution
  restriction.  The project includes $b=0$ as well as $b>0$: the first
  branch of the printed proposition allows $b=0$, whereas the printed
  corollary excludes only $b>0$.
\end{definition}

\begin{definition}[Finite page-extension stretching law]
  \label{def:finite-page-extension-stretching-law}
  \uses{def:finite-page-extension-nonliftable-crossing,def:page-extension}
  \lean{KIP126.Kervaire.Route.Tools.FinitePageExtensionStretchingLaw}
  \leanok
  A normalized family $P$ satisfies the finite stretching law if, for every
  \[
    2\le r'\le r<\infty,\qquad e(f)\le n\le r'-2+e(f),
  \]
  and actual later-page labels
  \[
    x\in Z_{r-1}^{s,t}(X),\qquad
    y\in Z_{r-1-n+e(f)}^{s+n,t+n}(Y),
  \]
  the earlier relation $d_n^{f,E_{r'}}(x)=y$ and the absence of a candidate
  from Definition~\ref{def:finite-page-extension-nonliftable-crossing}
  imply $d_n^{f,E_r}(x)=y$.
  This is the precise proposition to deliver for the chosen model; defining
  it does not prove that an arbitrary normalized family satisfies it.
  Its later-page target hypothesis and its exclusion of $b=0$ candidates
  make the project statement more restrictive than printed Corollary 6.21.
  The conclusion asserts existence of a later relation, without prescribing
  the reduction of every later representative to a specified earlier pair.
\end{definition}

% SOURCE: MainPaper/main.tex:1981-2068,
% Proposition prop:dec738d3; necessary-candidate direction only, with the
% shorter-page interpretation from Proposition prop:cross-f-Er.
\begin{proposition}[Certificate for loss of essentiality on an earlier page]
  \label{prop:page-extension-loss-certificate}
  \uses{prop:page-extension-restrict,def:finite-page-extension-nonliftable-crossing,prop:page-crossing-iff-synthetic-crossing}
  \notready
  For the normalized family of the chosen synthetic model, suppose an
  essential $(f,E_r)$ extension of length $n$, with source degree $(s,t)$,
  becomes inessential on $E_{r'}$, where
  \[
    2\le r'<r<\infty,\qquad e(f)\le n\le r'-2+e(f).
  \]
  Then there is a nonliftable crossing candidate from
  Definition~\ref{def:finite-page-extension-nonliftable-crossing}.
  This is only a necessary-candidate conclusion.  An arbitrary candidate
  in those degrees is not asserted to cause loss of essentiality of the
  chosen original extension.
\end{proposition}

\begin{proof}
  The planned proof follows the printed finite descent through essential
  shorter extensions.  The first branch gives a nonliftable source with
  $b=0$; later branches may give $b>0$.  The $\lambda$ comparison converts
  the essential contribution in the fixed quotient to an actual extension
  on page $r'-a$.  This conversion must retain the larger ordinary
  boundary exclusion in
  Proposition~\ref{prop:page-crossing-iff-synthetic-crossing}.
  The model comparisons and finite descent still require canonical Lean
  proofs.  No condition identifying the arbitrary candidate target with
  the original target has been added to the candidate definition.
\end{proof}

% SOURCE: MainPaper/main.tex:1981-2081,
% Proposition prop:dec738d3 and Corollary cor:dfc6043e; project finite,
% later-typed formulation with the stronger b>=0 crossing exclusion.
\begin{lemma}[First obstruction to a later page relation]
  \label{lem:page-extension-first-obstruction}
  \uses{def:finite-page-extension-nonliftable-crossing,def:finite-page-extension-stretching-law,prop:page-extension-loss-certificate,prop:page-crossing-iff-synthetic-crossing}
  \notready
  For the normalized family of the chosen synthetic model, take the ranges
  and later-page labels of
  Definition~\ref{def:finite-page-extension-stretching-law}.  If
  $d_n^{f,E_{r'}}(x)=y$ holds but $d_n^{f,E_r}(x)=y$ does not, then there
  is a candidate from
  Definition~\ref{def:finite-page-extension-nonliftable-crossing}.
  This is a relation-level claim.  It neither asserts that every failure
  to lift a specified strict solution produces such a candidate nor
  identifies the strict-solution cokernel with the candidate list.
\end{lemma}

\begin{proof}
  The required argument is the finite descent through shorter essential
  extensions in the actual $\widehat f$-ESS, using the $\rho$ and
  $\lambda$ comparison maps.  The comparison with a shorter page must
  preserve the larger ordinary boundary exclusion, as in
  Proposition~\ref{prop:page-crossing-iff-synthetic-crossing}.
  The canonical proof of this relation-level implication remains to be
  supplied.  Proposition 6.20 discusses loss of essentiality of a known
  later relation, and the printed Corollary 6.21 states existence of a
  later relation; their passage must be checked with the stated ranges.
  The generic cokernel criterion does not prove this step, and no
  surjectivity assertion about raw strict representatives is used here.
\end{proof}

% SOURCE: MainPaper/main.tex:2074-2081,
% Corollary cor:dfc6043e.  The infinite case additionally requires the
% project coherent-tower and untruncated comparison obligations below.
\begin{corollary}[Page-extension stretching]
  \label{cor:page-extension-stretch}
  \uses{def:finite-page-extension-stretching-law,lem:page-extension-first-obstruction,def:synthetic-ess-solution-tower,lem:synthetic-ess-coherent-tower-limit}
  \notready
  The normalized family of the chosen synthetic model satisfies the finite
  stretching law of
  Definition~\ref{def:finite-page-extension-stretching-law}.
  For permanent labels $x,y$, an earlier extension in the same range
  therefore yields the corresponding relation on every later finite page
  if all the nonliftable crossing candidates are excluded.

  For the untruncated conclusion $d_n^{f,E_\infty}(x)=y$, the finite
  strict solutions must additionally admit a coherent tower and the
  actual untruncated comparison of
  Lemma~\ref{lem:synthetic-ess-coherent-tower-limit} must be established
  for the model.  Explicit surjectivity of the actual adjacent restriction
  maps is one sufficient condition for constructing such a tower from an
  initial solution.  Alternatively, finiteness of every actual strict
  fiber, together with their nonemptiness, gives some coherent tower by
  Proposition~\ref{prop:synthetic-ess-coherent-tower-of-finite}.
  Exclusion of the finite candidates is not asserted to imply either
  surjectivity or finiteness of those representative sets.

  The two root predicates now fix the finite candidate list and the
  finite law, but the law's proof and its model comparison witnesses
  remain outstanding.  The untruncated comparison is a further proof
  obligation.  The retired free-operations statement expressed none of
  these conclusions.
\end{corollary}

\begin{proof}
  The finite law is the contrapositive of the first-obstruction lemma,
  whose model proof remains open.  The resulting finite relations give
  nonempty strict solution fibers.  To pass to $E_\infty$, construct a
  coherent tower using a separately justified compatibility argument and
  then prove that its limit gives an untruncated solution for $\widehat f$.
  Nonemptiness of each finite fiber alone completes neither step.
\end{proof}

% SOURCE: MainPaper/main.tex:2084-2099,
% Example exam:stretchext.  The printed example does not assert
% surjectivity of restrictions on all raw strict representative pairs.
\begin{proposition}[Stretching the Hopf-map extension to $E_\infty$]
  \label{prop:page-stretch-regression}
  \uses{prop:hopf-e3-page-extension,evidence:hopf-crossing-obstruction-exclusion,cor:page-extension-stretch,def:synthetic-ess-solution-tower,lem:synthetic-ess-coherent-tower-limit}
  \notready
  The $E_3$ extension
  $d_2^{\nu_{\mathrm{Hopf}},E_3}(h_0h_4^2)=h_0p$ stretches to the
  essential extension
  $d_2^{\nu_{\mathrm{Hopf}},E_\infty}(h_0h_4^2)=h_0p$.
  Its finite crossing exclusion supplies no assertion that every specified
  strict representative pair lifts between adjacent quotients.
\end{proposition}

\begin{proof}
  The independent finite Hopf certificate excludes the degree-allowed
  crossing candidates.  It also supplies that $h_0h_4^2$ and $h_0p$ are
  permanent representatives, giving both later-page label types.
  Once the finite stretching law is proved, these facts give an extension
  on every finite page.  To obtain the displayed untruncated conclusion,
  a coherent tower must still be constructed and compared to an actual
  untruncated solution, or an independent untruncated lifting argument
  must be supplied.  The finite candidate exclusion alone does not fill
  this gap.  Finally, the certificate's nonzero complete target-coset
  clause makes the resulting length-two relation essential, once its
  existence is established.
\end{proof}

% SOURCE: MainPaper/main.tex:1398-1405 and 929-961; project packaging of the
% zero-homology rotation needed by the normalized-cofiber interface.
\begin{lemma}[Normalized Hopf cofiber exactness case]
  \label{lem:normalized-hopf-exactness-case}
  \uses{evidence:normalized-hopf-detection,def:normalized-synthetic-exactness-case,prop:positive-adams-filtration-homology-zero,def:cofiber-triangle,def:hf2-homology,thm:external-synthetic-triangle-lift}
  \notready
  The classical Hopf cofiber triangle
  \[
    S^3\xrightarrow{\nu_{\mathrm{Hopf}}}S^0\xrightarrow g
    S^0/\nu_{\mathrm{Hopf}}\xrightarrow hS^4
  \]
  carries the zero constructor of
  Definition~\ref{def:normalized-synthetic-exactness-case}.  Explicitly,
  \[
    H_*\nu_{\mathrm{Hopf}}=0,
    \qquad
    (e(\nu_{\mathrm{Hopf}}),e(g),e(h))=(1,0,0),
  \]
  and its long exact homology sequence supplies the rotated short exact
  sequence.  The constructor's lifted-triangle field is exactly the witness
  retained by Proposition~\ref{evidence:normalized-hopf-detection}, and its
  normalized Hopf component is therefore equal to $[h_2]$.
\end{lemma}

\begin{proof}
  The filtration-one detection by $h_2$ gives
  $e(\nu_{\mathrm{Hopf}})=1$, while degree considerations make the Hopf map
  zero on $H\mathbb F_2$-homology.  In the cofiber long exact sequence, $g$
  is the nonzero bottom-cell inclusion on homology and $h$ is the nonzero
  top-cell quotient.  The contrapositive of
  Proposition~\ref{prop:positive-adams-filtration-homology-zero} therefore
  gives Adams filtration zero and $e(g)=e(h)=0$.  The remaining
  long-exact-sequence terms split by degree into
  the required short exact sequence.  Package these facts and the displayed
  $e$-pattern around the precise lifted-triangle witness retained by the
  normalized-detection evidence; do not invoke the nonunique lift theorem a
  second time.  The evidence's component equality becomes the zero tag's
  equality $\widehat{\nu_{\mathrm{Hopf}}}=[h_2]$.
\end{proof}
